%======================================================================
\section{Proof of the Main Theorem}\label{sec:proof}

We follow the two steps described in the introduction.
Section~\ref{sec:criterion} proves the criterion (step~1).
Section~\ref{sec:decoding} proves the decoding (step~2) and
Theorem~\ref{thm:main}.

\subsection{A criterion for minimality}\label{sec:criterion}

\begin{definition}\label{def:onestep}
Let $\delta$ be a formula, let $Q$ be a nonempty set of occurrences of
$\delta$ in $\alpha$ (that is, $\alpha_u=\delta$ for $u\in Q$), and let $z$
be a variable not occurring in $\alpha$. If $\delta$ is a variable, assume
moreover that some occurrence of $\delta$ in $\alpha$ is not in $Q$. Then
$\alpha[Q:=z]$ is called a \emph{one-step generalization} of $\alpha$.
\end{definition}

\begin{example}
For $\alpha=(a\to b)\to(a\to b)$, both $z\to(a\to b)$ (with $\delta=a\to b$)
and $(z\to b)\to(a\to b)$ (with $\delta=a$) are one-step generalizations of
$\alpha$. $(z\to b)\to(z\to b)$ is not one, because all occurrences of $a$
are replaced.
\end{example}

\begin{lemma}\label{lem:strict}
Every one-step generalization $\beta=\alpha[Q:=z]$ of $\alpha$ satisfies
$\beta\ssub\alpha$.
\end{lemma}
\begin{proof}
$\beta\sub\alpha$, because substituting $\delta$ for $z$ in $\beta$ gives
$\alpha$ ($z$ does not occur in $\alpha$). Suppose, for a contradiction, that
$\tau\alpha=\beta$ for some $\tau$. Call connectives and constants
\emph{symbols}. A substitution never decreases the number of occurrences of
symbols.

If $\delta$ is not a variable, it contains a symbol, so $\beta$ has fewer
symbols than $\alpha$. This is a contradiction.

If $\delta$ is a variable $p$, then $\alpha$ and $\beta$ have the same number
of symbols, so $\tau$ maps each variable of $\alpha$ to a variable. Then
$\tau\alpha$ has the same syntax tree as $\alpha$, with the variable
$\tau(p)$ at every leaf where $\alpha$ has $p$. But at these leaves $\beta$
has $z$ (at $u\in Q$) and $p$ (at the occurrence outside $Q$). Hence
$\tau(p)=z$ and $\tau(p)=p$, a contradiction.
\end{proof}

\begin{lemma}\label{lem:criterion}
Let $L$ be a logic and $\alpha\in L$. If no one-step generalization of
$\alpha$ is in $L$, then $\alpha$ is minimal in $L$.
\end{lemma}
\begin{proof}
Let $\gamma\in L$ and $\sigma\gamma=\alpha$. We show $\alpha\sub\gamma$.

Suppose first that $\sigma$ maps the variables of $\gamma$ injectively to
variables. Then there is a substitution $\tau$ with $\tau(\sigma(q))=q$ for
every variable $q$ of $\gamma$, and $\tau\alpha=\tau\sigma\gamma=\gamma$.

Otherwise, there is a variable $q$ of $\gamma$ such that either (i)
$\sigma(q)$ is not a variable, or (ii) $\sigma(q)=\sigma(q')$ is a variable
for some variable $q'\ne q$ of $\gamma$. Let $z$ be a variable not occurring
in $\alpha$, and let $\sigma'$ agree with $\sigma$ except that
$\sigma'(q)=z$. Every occurrence of $q$ in $\gamma$ yields an occurrence of
$\delta:=\sigma(q)$ in $\alpha=\sigma\gamma$. Let $Q$ be the set of these
occurrences; $Q\ne\emptyset$. Then $\sigma'\gamma=\alpha[Q:=z]$. In case (ii),
the occurrences of $q'$ in $\gamma$ yield occurrences of the variable
$\delta$ that are not in $Q$. So $\sigma'\gamma$ is a one-step generalization
of $\alpha$, and $\sigma'\gamma\in L$ by (C1). This contradicts the
hypothesis.
\end{proof}

\begin{corollary}\label{cor:criterion}
Let $L$ be a logic. A formula $\alpha\in L$ is minimal in $L$ if and only if
no one-step generalization of $\alpha$ is in $L$.
\end{corollary}
\begin{proof}
``If'' is Lemma~\ref{lem:criterion}. ``Only if'' follows from
Lemma~\ref{lem:strict}.
\end{proof}


\subsection{Decoding one-step generalizations}\label{sec:decoding}

For an occurrence $u$ of $\alpha$, the \emph{reading} of $u$ is the
occurrence of $r_u$ inside the body $b_v$ of its parent $v$, or the last
formula $r_\varepsilon$ of $\St(\alpha)$ when $u=\varepsilon$. The next lemma
lists the facts about the shape of $\St(\alpha)$ that we use. All of them are
immediate from the definition.

\begin{lemma}\label{lem:shape}
\begin{enumerate}
\item[(a)] For each compound occurrence $u$, $x_u$ occurs in $\St(\alpha)$
  exactly twice: at the reading of $u$, and as the \emph{head} of $E_u$
  (the variable side of the link).
\item[(b)] The occurrences of atoms of $\alpha$ in $\St(\alpha)$ are exactly
  the readings of the leaves of $\alpha$. So for each atom $c$ they
  correspond bijectively to the leaves $u$ with $\alpha_u=c$.
\item[(c)] Every compound subformula occurrence in $\St(\alpha)$ is of
  exactly one of the following kinds:
  \begin{itemize}
  \item a body $b_u$ (both immediate subformulas are atoms);
  \item a link $E_u$ (one immediate subformula is an atom, the other a body);
  \item a \emph{tail} $T_i=E_{u_i}\to\dots\to E_{u_m}\to r_\varepsilon$ with
    $1\le i\le m$ (its left immediate subformula is a link).
  \end{itemize}
  In particular, no tail is a subformula of a link, and distinct tails are
  distinct formulas.
\item[(d)] If $b_u=b_v$ as formulas, then $\alpha_u=\alpha_v$. If $E_u=E_v$
  as formulas, then $u=v$.
\end{enumerate}
\end{lemma}
\begin{proof}
(a)--(c) follow by inspecting Definition~\ref{def:S}. For the last claims
of (c): the proper subformulas of a link are its body and atoms, and $T_i$
has $m-i+1$ premises. For (d), note that $r_w$ determines $\alpha_w$: if
$r_w$ is an atom of $\alpha$, then $\alpha_w=r_w$; if $r_w=x_{w'}$, then
$w=w'$. So $b_u=b_v$ gives the same main connective,
$\alpha_{u0}=\alpha_{v0}$ and $\alpha_{u1}=\alpha_{v1}$, hence
$\alpha_u=\alpha_v$. A positive link has the form (body)$\to$(variable), and
a negative one (variable)$\to$(body). So $E_u=E_v$ forces equal polarity and
$x_u=x_v$, hence $u=v$.
\end{proof}

We call a one-step generalization $\St(\alpha)[P:=z]$ \emph{normal} if every
position in $P$ is a reading or a body. We first reduce to normal ones.
Throughout, $z$ is a variable not occurring in $\St(\alpha)$.

\begin{lemma}[Normalization]\label{lem:normal}
If some one-step generalization of $\St(\alpha)$ belongs to a logic $L$,
then some normal one-step generalization of $\St(\alpha)$ belongs to $L$.
\end{lemma}
\begin{proof}
Let $\alpha'=\St(\alpha)[P:=z]\in L$, where $P$ is a set of occurrences of
$\delta$. We go through the kinds of $\delta$ given by
Lemma~\ref{lem:shape}. Each step replaces $\alpha'$ by a substitution
instance, which stays in $L$ by (C1).
\begin{itemize}
\item \emph{$\delta$ is an atom of $\alpha$ or a body.} Then every
  occurrence of $\delta$ is a reading, resp.\ a body, by
  Lemma~\ref{lem:shape}(b),(c). So $\alpha'$ is already normal.
\item \emph{$\delta=x_u$.} By Lemma~\ref{lem:shape}(a) and the definition of
  a one-step generalization, $P$ is one of the two occurrences of $x_u$. If
  it is the reading, $\alpha'$ is normal. If it is the head, exchange the
  variables $x_u$ and $z$. The result is $\St(\alpha)$ with the reading of
  $u$ replaced by $z$, which is normal.
\item \emph{$\delta=E_u$.} By Lemma~\ref{lem:shape}(c),(d), $E_u$ occurs
  only once, so $\alpha'$ is $\St(\alpha)$ with the premise $E_u$ replaced
  by $z$. Substitute for $z$ the link $E_u$ with its head replaced by $z$.
  This gives $\St(\alpha)$ with the head of $E_u$ replaced by $z$, which is
  the previous case.
\item \emph{$\delta=T_i$.} By Lemma~\ref{lem:shape}(c), $T_i$ occurs only
  once, so $\alpha'=E_{u_1}\to\dots\to E_{u_{i-1}}\to z$. Substitute for $z$
  the tail $T_i$ with its last formula $r_\varepsilon$ replaced by $z$. This
  gives $\St(\alpha)$ with the reading of $\varepsilon$ replaced by $z$,
  which is normal.\qedhere
\end{itemize}
\end{proof}

\begin{lemma}[Decoding]\label{lem:key}
Let $\alpha'=\St(\alpha)[P:=z]$ be a normal one-step generalization. Let $Q$
be the set of occurrences $u$ of $\alpha$ such that the reading or the body
of $u$ lies in $P$. Then $\beta=\alpha[Q:=z]$ is a one-step generalization of
$\alpha$, and there is a substitution $\theta$ with
\[
\theta\alpha'=(\gamma_1\to\gamma_1)\to\dots\to(\gamma_k\to\gamma_k)\to\beta .
\]
\end{lemma}
\begin{proof}
\emph{$\beta$ is a one-step generalization.} $Q$ is nonempty. All
$\alpha_u$ ($u\in Q$) are the same formula $\delta_0$:
\begin{itemize}
\item if $\delta$ is an atom, then $\delta_0=\delta$;
\item if $\delta=x_w$, then $Q=\{w\}$;
\item if $\delta$ is a body, use Lemma~\ref{lem:shape}(d).
\end{itemize}
In particular, no element of $Q$ is strictly below another. If $\delta_0$ is
a variable, then $\delta=\delta_0$ and $Q$ misses some leaf labelled
$\delta_0$ by Lemma~\ref{lem:shape}(b).

\emph{The substitution.} Call $v$ \emph{hidden} if it is strictly below some
element of $Q$. For $v$ not hidden, write $\beta_v$ for the subformula of
$\beta$ at $v$ ($\beta_v=z$ for $v\in Q$). Let $\theta$ fix $z$ and the
variables of $\alpha$, and put
\[
\theta(x_v)=\begin{cases}
\alpha_v,&\text{if $v$ is hidden, or the reading of $v$ is in $P$},\\
\beta_v,&\text{otherwise.}
\end{cases}
\]
Here is the rule behind this choice: $x_v$ is sent to $\alpha_v$ when only
its head survives, and to $\beta_v$ when its reading survives.

\emph{Readings.} If $v$ is hidden, $\theta$ maps the reading of $v$ to
$\alpha_v$. If $v$ is not hidden, it maps it to $\beta_v$:
\begin{itemize}
\item if the reading is in $P$, it is $z=\beta_v$;
\item otherwise it is $r_v$, and $\theta(r_v)=\beta_v$ by definition (for a
  leaf $v\notin Q$, $\beta_v=\alpha_v=r_v$).
\end{itemize}
(Positions below an element of $Q$ are not in $P$, because $Q$ has no
nested elements.)

\emph{Premises.} Let $v$ be compound, with main connective $\circ$.
\begin{itemize}
\item If $v$ is hidden, $E_v$ is unchanged, and $\theta$ maps its body to
  $\alpha_{v0}\circ\alpha_{v1}=\alpha_v=\theta(x_v)$.
\item If $v$ is not hidden and $v\notin Q$, $E_v$ is unchanged, and $\theta$
  maps its body to $\beta_{v0}\circ\beta_{v1}=\beta_v=\theta(x_v)$.
\item If $v\in Q$ and its body is in $P$, then $E_v$ has become $z\to x_v$ or
  $x_v\to z$, and $\theta(x_v)=\beta_v=z$.
\item If $v\in Q$ and its reading is in $P$, then $E_v$ is unchanged. Its
  children are hidden, so $\theta$ maps its body to $\alpha_v=\theta(x_v)$.
\end{itemize}
In every case the premise is mapped to $\gamma\to\gamma$ (in either
polarity).

\emph{Conclusion.} The last formula is the reading of $\varepsilon$, which is
not hidden. So it is mapped to $\beta_\varepsilon=\beta$.
\end{proof}

\begin{proof}[Proof of Theorem~\ref{thm:main}]
If $\alpha$ is an atom, then $\St(\alpha)=\alpha$ and there is nothing to
prove. Otherwise, suppose some one-step generalization of $\St(\alpha)$ is in
$L$. By Lemma~\ref{lem:normal}, some normal one-step generalization $\alpha'$
is in $L$. By Lemma~\ref{lem:key} and (C1)--(C3), the one-step
generalization $\beta$ of $\alpha$ is in $L$. This contradicts the
minimality of $\alpha$ (Corollary~\ref{cor:criterion}). Hence no one-step
generalization of $\St(\alpha)$ is in $L$. Since $\St(\alpha)\in L$,
$\St(\alpha)$ is minimal by Corollary~\ref{cor:criterion}.
\end{proof}

Consistency of $L$ is not needed. If $\alpha$ is compound and minimal in $L$,
then the variable $z$, a one-step generalization of $\alpha$, is not in $L$.

\begin{example}\label{ex:decode}
We illustrate Lemma~\ref{lem:key} with Example~\ref{ex:peirce}. Replace the
reading of $00$, i.e.\ the second occurrence of $w=x_{00}$, by $z$:
\[
\alpha'=((a\to b)\to w)\to(y\to(z\to a))\to((y\to a)\to x)\to x .
\]
Here $Q=\{00\}$ and $\beta=(z\to a)\to a$. The substitution is
$w\mapsto a\to b$ (only the head of $w$ survives), $y\mapsto z\to a$ and
$x\mapsto\beta$:
\[
\theta\alpha'=((a\to b)\to(a\to b))\to((z\to a)\to(z\to a))\to
(\beta\to\beta)\to\beta .
\]
Since $\beta\notin\CL$ ($z=0$, $a=0$), neither is $\alpha'$.
\end{example}
